\documentclass[aspectratio=169]{beamer}

\usetheme{Madrid}
\usecolortheme{default}

\usepackage{listings}
\usepackage{booktabs}
\usepackage{hyperref}
\usepackage{xcolor}
\usepackage{array}
\usepackage{graphicx}

% Code listing style
\lstset{
    basicstyle=\ttfamily\small,
    breaklines=true,
    backgroundcolor=\color{gray!10},
    keywordstyle=\color{blue},
    commentstyle=\color{green!50!black},
    stringstyle=\color{red!70!black},
    showstringspaces=false
}

\title{Scientific AI}
\subtitle{D-Lab AI Pulse Workshop}
\author{}
\date{}

\begin{document}

\begin{frame}[plain]
    \titlepage
\end{frame}

% =============================================================================
% INTRODUCTION
% =============================================================================

\begin{frame}{The Promise and the Problem}
    \textbf{There has been a lot of interest in using AI tools for academic research.}

    \vspace{1em}
    \textbf{What they're great at:}
    \begin{itemize}
        \item Brainstorming and drafting
        \item Explaining complex concepts
        \item Writing and editing code
        \item Summarizing long documents
    \end{itemize}

    \vspace{1em}
    \textbf{But for academic research, there's a critical flaw:}

    \vspace{0.5em}
    \textit{These tools can present fabricated information as fact and without warning.}

    \vspace{1em}
    Today: Why this happens, and tools designed specifically to avoid it.
\end{frame}

\begin{frame}{Today's Roadmap}
    \textbf{1. The Problem}
    \begin{itemize}
        \item Why general AI hallucinates citations
        \item Real-world consequences
    \end{itemize}

    \vspace{1em}
    \textbf{2. The Solution}
    \begin{itemize}
        \item How specialized research tools work differently
        \item Perplexity, Consensus, Elicit
    \end{itemize}

    \vspace{1em}
    \textbf{3. The Frontier}
    \begin{itemize}
        \item Kosmos: Autonomous research agents
        \item What's coming next
    \end{itemize}

    \vspace{1em}
    \textbf{Live demos throughout.}
\end{frame}

% =============================================================================
% PART 1: THE PROBLEM
% =============================================================================

\begin{frame}{The Hallucination Problem}
    \textbf{Large Language Models predict plausible text, not accurate text.}

    \vspace{0.5em}
    They are not designed to verify whether their output is true, and lack uncertainty quantification.

    \vspace{1em}
    \textbf{Why it's convincing:} A fake citation like ``Smith \& Johnson (2020)'' follows patterns from millions of real papers --- common surnames, recent year, proper format. The model generates \textit{plausible} text, not \textit{verified} text.

    \vspace{1em}
  \textbf{The numbers:}
  \begin{itemize}
      \item ChatGPT fabricates \textbf{18--55\%} of citations depending on model version\textsuperscript{1}
      \item \textbf{24--43\%} of real references contain substantive errors\textsuperscript{1}
      \item Specialized legal AI tools hallucinate in \textbf{17--33\%} of queries\textsuperscript{2}
  \end{itemize}

  \vspace{0.3em}
  {\scriptsize \textsuperscript{1}Walters \& Wilder (2023), \textit{Scientific Reports};
  \textsuperscript{2}Magesh et al.\ (2025), \textit{J.\ Empirical Legal Studies}}
\end{frame}

  \begin{frame}{Real-World Consequences}
      \textbf{In consulting:}
      \begin{itemize}
          \item Deloitte Australia (2025): Welfare report with fabricated references
          \item Deloitte Canada (2025): Healthcare report with fake citations
      \end{itemize}

      \vspace{0.5em}
      \textbf{In law:}
      \begin{itemize}
          \item Mata v.\ Avianca (2023): 6 nonexistent cases cited
          \item California appeal (2025): 21 of 23 quoted cases fabricated
          \item Hundreds of errant legal citations flagged in US courts to date
      \end{itemize}

      \vspace{0.5em}
      \textbf{In academia:}
      \begin{itemize}
          \item 100+ hallucinated citations found across 51 accepted NeurIPS 2025 papers
          \item Fake citations get published, cited by others, and the problem cascades
      \end{itemize}
  \end{frame}

% =============================================================================
% PART 2: THE SOLUTION
% =============================================================================

  \begin{frame}{Why Specialized Tools?}
      \begin{table}
          \centering
          \begin{tabular}{p{6cm}p{6cm}}
              \toprule
              \textbf{General AI} & \textbf{Academic AI} \\
              \footnotesize (ChatGPT, Claude, Gemini) & \footnotesize (Elicit, Consensus, Perplexity) \\
              \midrule
              Generates plausible text & Retrieves from verified sources \\
              No database verification & 200M+ papers (Semantic Scholar, OpenAlex, PubMed) \\
              Citations may be fabricated & Citations are real and verifiable \\
              \bottomrule
          \end{tabular}
      \end{table}

      \vspace{1em}
      \textbf{Key difference:} \textit{Retrieval} from verified sources vs.\ \textit{generation} from patterns.
  \end{frame}

  \begin{frame}{Three Tools, One Workflow}
      \textbf{All three:} Ask a question $\rightarrow$ get papers with AI summaries $\rightarrow$ filter $\rightarrow$ export
  citations

      \vspace{0.5em}
      \textbf{All three:} Semantic search, sentence-level citations, citation export

      \vspace{1em}
      \begin{table}
          \centering
          \begin{tabular}{lll}
              \toprule
              \textbf{Tool} & \textbf{Best For} & \textbf{Unique Feature} \\
              \midrule
              Perplexity & Quick context, broad questions & Web + academic search \\
              Consensus & ``What does evidence say?'' & Consensus Meter \\
              Elicit & Systematic reviews & Structured data extraction \\
              \bottomrule
          \end{tabular}
      \end{table}

      \vspace{1em}
      \textbf{Limitations:} Smaller coverage than Google Scholar or library databases. Better for established topics. Supplement,
  don't replace, traditional database searches.
  \end{frame}

% =============================================================================
% PART 3: TOOL DEEP DIVES
% =============================================================================

\begin{frame}[fragile]{Perplexity}
    \textbf{Broad search:} Academic papers + web sources + news

    \vspace{1em}
    \textbf{Features:}
    \begin{itemize}
        \item \textbf{Model selection:} Choose between Claude, GPT, Grok --- different strengths per query
        \item \textbf{Academic Mode:} Restricts sources to peer-reviewed literature
        \item \textbf{Deep Research:} Breaks question into sub-questions, synthesizes across dozens of sources
        \item \textbf{Model Council:} Queries multiple models and synthesizes their answers
    \end{itemize}

    \vspace{1em}
    \textbf{Best for:} Background research, exploring unfamiliar topics, quick fact-checking

    \vspace{1em}
    \textbf{Demo:}
    \begin{lstlisting}[basicstyle=\ttfamily\small]
What are the main criticisms of DSGE models
in macroeconomics?
    \end{lstlisting}
\end{frame}

\begin{frame}{Perplexity Spaces}
    \textbf{Persistent, shareable research environments.}

    \vspace{0.5em}
    \textbf{What you can do:}
    \begin{itemize}
        \item \textbf{Upload files} --- syllabi, readings, datasets, codebooks, notes
        \item \textbf{Invite collaborators} --- classmates, lab members, advisors
        \item \textbf{Ask questions} grounded in your uploaded materials + the full web
        \item \textbf{Build a shared research trail} everyone can see and build on
    \end{itemize}

    \vspace{0.5em}
    \textbf{Example workflow:}

    \vspace{0.3em}
    \textit{Create a Space for your final project $\rightarrow$ upload your readings and dataset $\rightarrow$ invite your partner $\rightarrow$ now you both ask questions and build a shared evidence base.}
\end{frame}

\begin{frame}[fragile]{Consensus}
    \textbf{Focused on academic literature:} Peer-reviewed papers, preprints, and conference proceedings

    \vspace{0.5em}
    \textbf{Includes:} Wiley, Taylor \& Francis, and other major publishers

    \vspace{0.5em}
    \textbf{Features:} Consensus Meter, import from Zotero, upload your own papers

    \vspace{0.5em}
    Premium available for Berkeley

    \vspace{1em}
    \textbf{Best for:} Evidence-based questions, checking scientific consensus, policy debates

    \vspace{1em}
    \textbf{Demo:}
    \begin{lstlisting}[basicstyle=\ttfamily\small]
Does minimum wage reduce employment?
    \end{lstlisting}
\end{frame}

\begin{frame}[fragile]{Elicit}
    \textbf{Designed for systematic reviews:} Screening, data extraction, report generation

    \vspace{1em}
    \textbf{Workflows:}
    \begin{itemize}
        \item \textbf{Extract Data} --- build custom tables from papers (sample sizes, effect sizes, methods)
        \item \textbf{Competitive Landscape} --- compare products or approaches across academic and non-academic sources
        \item \textbf{General Research Agent} --- flexible agent that breaks down any research question
    \end{itemize}

    \vspace{1em}
    \textbf{Demo:}
    \begin{lstlisting}[basicstyle=\ttfamily\small]
What methods are used to measure inflation
expectations?
    \end{lstlisting}
\end{frame}

\begin{frame}[fragile]{JSTOR AI Research Tool}
    \textbf{AI built on top of a library database you already have access to.}

    \vspace{1em}
    \textbf{What it does:}
    \begin{itemize}
        \item Assess content relevance --- surfaces key points and arguments from texts
        \item Discover related topics and materials across the JSTOR corpus
        \item Conversational --- use natural language to search or ask questions about what you're reading
    \end{itemize}

    \vspace{1em}
    \textbf{Demo:}
    \begin{lstlisting}[basicstyle=\ttfamily\small]
What is the relationship between central bank
independence and inflation?
    \end{lstlisting}
\end{frame}

\begin{frame}{When to Use Something Else}
    \textbf{These tools are slower than ChatGPT.} That's by design --- more verification takes more time. Use ChatGPT/Claude for brainstorming and drafting; use these tools for sourcing.

    \vspace{0.5em}
    \textbf{Coverage gaps:}
    \begin{itemize}
        \item Very recent papers ($<$3 months) --- indexing takes time
        \item Niche topics with thin literature
        \item Non-English publications
    \end{itemize}

    \vspace{0.5em}
    \textbf{Use Google Scholar or your library database when you need:}
    \begin{itemize}
        \item Exhaustive coverage --- dissertations, theses, grey literature, working papers
        \item Known-item searches --- you already have a title or DOI
        \item Full-text access through Berkeley's subscriptions
        \item Non-English or region-specific literature
    \end{itemize}

    \vspace{0.5em}
    \textbf{Best approach:} Start broad with these tools, verify and fill gaps with traditional databases.
\end{frame}

\begin{frame}{Discussion}
    \begin{enumerate}
        \item \textbf{Teach or restrict?} Students and researchers will use these tools regardless. Are we better off trying to stop them, or teaching them to use them well?

        \item \textbf{Industry is pushing AI.} Not just welcoming it, but actively requiring it. How does that change what and how we teach?

        \item \textbf{The next shift.} Information went from low volume, high signal-to-noise, to high volume, low signal-to-noise, to on-demand volume and unobservable signal-to-noise ratio. We can't imagine research without Google today. What will research look like in 25 years?

        \item \textbf{Equity and access.} Free tiers exist, but premium features cost money. Does AI narrow the gap between well-funded and under-funded researchers, or widen it?
    \end{enumerate}
\end{frame}



% =============================================================================
% PART 4: JSTOR & KOSMOS
% =============================================================================



\begin{frame}{What is Kosmos?}
    \textbf{An autonomous research agent} from FutureHouse / Edison Scientific

    \vspace{0.3em}
    \begin{center}
        \includegraphics[width=0.65\textwidth]{../690ac9822b6bc833f1f47b8a_Screenshot_2025-11-04_at_11.32.50_AM.png}
    \end{center}

    \vspace{0.3em}
    \textbf{Two modes:}
    \begin{itemize}
        \item \textbf{Kosmos} (full automation): Submit a question + data (up to 15GB), receive a report in $\sim$12 hours
        \item \textbf{Interactive agents:} Literature synthesis, data analysis, molecular design, precedent search
    \end{itemize}

    \vspace{0.3em}
    \textbf{Access:} Free for academics with .edu email (650 credits/month)
\end{frame}

\begin{frame}{Kosmos: Discoveries}
    {\footnotesize Mitchener et al.\ (2025), arXiv:2511.02824 --- 7 discoveries across 4 fields.
    Estimated equivalent of $\sim$6 months of human research per run.}

    \vspace{0.2em}
    \begin{columns}[T]
    \begin{column}{0.35\textwidth}
        \includegraphics[width=\textwidth]{../690ac9c1d363c9893d2c868f_Screenshot_2025-11-04_at_11.33.15_AM.png}
    \end{column}
    \begin{column}{0.63\textwidth}
    {\footnotesize
    \textbf{Novel findings:}
    \begin{itemize}\setlength{\itemsep}{1pt}
        \item SOD2 causally reduces myocardial fibrosis (Mendelian randomization)
        \item GWAS variant regulates SSR1 in Type 2 diabetes (multi-omic analysis)
        \item Temporal ordering of ECM decline during tau accumulation in Alzheimer's
        \item Flippase downregulation triggers neuron loss in entorhinal cortex
    \end{itemize}

    \vspace{0.2em}
    \textbf{Independent reproductions:}
    \begin{itemize}\setlength{\itemsep}{1pt}
        \item Nucleotide metabolism in hypothermic brain protection
        \item Humidity as critical factor in perovskite solar cell fabrication
        \item Log-normal connectivity distributions across 5 species
    \end{itemize}
    }
    \end{column}
    \end{columns}
\end{frame}

\begin{frame}{Kosmos: Best Practices}
    \begin{enumerate}
        \item \textbf{Clear, feasible research objective.} Single, well-defined goal with scope for iterative testing. Should require weeks or months of human work, not a single analysis.
        \item \textbf{Provide enough context.} Phrase it as you would for an experienced colleague who just joined your team. Include field-specific nuances and experimental design choices.
        \item \textbf{Describe your data.} Label columns intuitively or provide a codebook. A collaborator should be able to interpret the data without further clarification.
        \item \textbf{Complex data works best.} High-dimensional data (scRNA-seq, proteomics, multi-sample environmental data). Multiple datasets welcome, up to 15GB total.
        \item \textbf{Processed data.} Kosmos can do QC, but works best on cleaned, processed data. Not recommended for raw sequencing or raw imaging.
        \item \textbf{Iterate.} Start with a familiar dataset. Runtime scales with data size, so begin small.
    \end{enumerate}

    \vspace{0.3em}
    {\footnotesize Full guidelines: \url{https://edisonscientific.com/articles/kosmos-best-practices-guidelines}}
\end{frame}

% =============================================================================
% PART 5: RESOURCES
% =============================================================================

\begin{frame}{Resources: Tools}
    {\small
    \textbf{Tools covered today:}
    \begin{itemize}
        \item Perplexity: \url{https://perplexity.ai/}
        \item Consensus: \url{https://consensus.app/}
        \item Elicit: \url{https://elicit.com/}
        \item JSTOR AI: \url{https://www.jstor.org/} (Free via Berkeley Library)
        \item Kosmos / Edison Scientific: \url{https://edisonscientific.com/}
    \end{itemize}

    \vspace{0.5em}
    \textbf{Also mentioned:}
    \begin{itemize}
        \item Research Rabbit: \url{https://www.researchrabbit.ai/}
        \item Connected Papers: \url{https://www.connectedpapers.com/}
        \item Litmaps: \url{https://www.litmaps.com/}
        \item Scite.ai: \url{https://scite.ai/}
        \item SciSpace: \url{https://typeset.io/}
    \end{itemize}
    }
\end{frame}

\begin{frame}{Resources: References}
    {\small
    \textbf{Papers and data cited:}
    \begin{itemize}
        \item Walters \& Wilder (2023). ``Fabrication and Errors in the Bibliographic Citations Generated by ChatGPT.'' \textit{Scientific Reports}. \url{https://doi.org/10.1038/s41598-023-41032-5}
        \item Magesh et al.\ (2025). ``Hallucination-Free? Assessing the Reliability of Leading AI Legal Research Tools.'' \textit{J.\ Empirical Legal Studies}.
        \item Ansari (2026). ``Compound Deception in Elite Peer Review: A Failure Mode Taxonomy of 100 Fabricated Citations at NeurIPS 2025.'' arXiv:2602.05930. \url{https://arxiv.org/abs/2602.05930}
        \item Mitchener et al.\ (2025). ``Kosmos: Autonomous Research Agent.'' arXiv:2511.02824.
        \item Charlotin, D. AI Hallucination Cases Database (912 cases). \url{https://www.damiencharlotin.com/hallucinations/}
    \end{itemize}

    \vspace{0.5em}
    \textbf{D-Lab:} \url{https://dlab.berkeley.edu/}
    }
\end{frame}

\end{document}
